% Real canonical-vs-delta KB examples (Tomato), from *_delta_validation.md.
\begin{table*}[t]
\centering\small
\setlength{\tabcolsep}{5pt}
\renewcommand{\arraystretch}{1.15}
\caption{\textbf{Canonical KB vs.\ image-grounded Delta KB} (Tomato). The
\emph{canonical} column is verbatim, source-cited extension text; the
\emph{delta} is what a VLM reports from a state-specific \bugwood{} field
photograph of that disease. \emph{Grounding} rows confirm the canonical claim
against a real field image; \emph{Refinement} rows show the delta extending,
re-staging, or adding look-alikes the canonical text omits---the signal the
inference re-ranker consumes.}
\label{tab:kb-delta}
\begin{tabular}{@{}llp{0.36\textwidth}p{0.34\textwidth}l@{}}
\toprule
Disease & Field & Canonical KB (cited) & Delta KB (\texttt{image\_shows}) & Type \\
\midrule
\multicolumn{5}{@{}l}{\textit{Grounding --- delta confirms canonical against a field image}}\\
Late Blight & fruit & Dark brown, greasy-to-leathery circular spots. & Tomato fruits with dark brown, greasy-to-leathery circular spots. & match\\
Anthracnose & sporulation & Masses of salmon-colored sausage-shaped spores may form on the lesion surface. & Salmon-colored sausage-shaped masses forming on the lesion surface. & match\\
Southern Blight & look-alikes & Tan to reddish-brown spherical sclerotia (1--2\,mm) at the stem base. & Tan to reddish-brown spherical sclerotia (1--2\,mm) at the stem base. & match\\
\midrule
\multicolumn{5}{@{}l}{\textit{Refinement --- delta extends / re-stages / adds look-alikes}}\\
Tomato Leaf Mould & lesion & Pale green / yellowish spots; olive-green mould on the lower leaf surface. & Brown spots clustered near major veins, some discrete. & stage/color\\
Bacterial Speck & look-alikes & Bacterial spot; spotted wilt; bacterial canker. & Target-spot lesion: central dark-brown area within concentric rings. & added look-alike\\
\bottomrule
\end{tabular}
\end{table*}
